% Source: https://artofproblemsolving.com/community/c3h2550626_seemingly_correct_solution_implies_easily_disprovable_corollary
% Stable identity: aops-topic-2550626-q01
A student ``proves'' that $n$ and $n^5$ have the same last digit for every positive integer $n$ as follows:

``If $n$ and $n^5$ have the same last digit, then $n^5\equiv n\pmod{10}$. Since $5$ is a factor of $10$, $n^5\equiv n\pmod5$, which is true by Fermat's theorem.''

The same argument with $2$ in place of $5$ seems to ``prove'' that $n$ and $n^2$ always have the same last digit.
\begin{enumerate}
\item[(a)] Explain what is wrong with the argument.
\item[(b)] Give a correct proof that $n$ and $n^5$ have the same last digit for every positive integer $n$.
\item[(c)] Show that $n$ and $n^2$ need not have the same last digit.
\end{enumerate}
